% Self-contained mathematical report fragment for the new AX11GDH command system.
% Requires amsmath, amssymb, amsthm, booktabs. No custom macros required.
\section{The AX11GDH command system and its iPRA preservation proof}
\label{sec:ax11gdh-preservation}

The new command system uses the same numerical calculus $\mathcal C_4$
as the ordinary AX11H parent. The change concerns which independent
command operands supply two emitted formulas. The diagonal entry is
$\Gamma(D)$, with $D$ the accumulated offset, while the generated
universal formula is $U(c,h)$, with $c$ the explicit AX tag parameter.
These two choices require different diagonal and generalization selector
words. The following proof supplies those words explicitly and retains
the real fourth AX operand and the cached-only halt test.

\subsection{Exact equations and all eleven stack entries}

Put
\[
 p(a,b)=\frac{(a+b)(a+b+1)}2+b,\quad
 I(A,B)=4p(A,B),\quad M(a,b)=4p(a,b)+1,
\]
\[
 U(a,B)=4p(B,a)+2,\qquad E(a,b)=4p(a,b)+3,
 \qquad N(B)=I(B,1).
\]
Conjunction is $N(I(A,N(B)))$, and
$\exists_a B=N(U(a,N(B)))$. Zero is truth and one is the designated
contradiction. The ordinary six heads are
\begin{align*}
 h_1(D,H)&=I(\exists_D U(D,H),H),\\
 h_2(D,A,H)&=I(U(D,U(A,H)),U(A,U(D,H))),\\
 h_3(D,A,H)&=I(M(H,D),I(E(H,A),M(A,D))),\\
 h_4(D,A,H)&=I(E(H,D),I(M(A,H),M(A,D))),\\
 h_5(D,H)&=\exists_D E(H,D),\\
 h_6(D,A,H)&=I(U(D,A),I(U(D,I(A,H)),U(D,H))).
\end{align*}
The guarded gate is defined by
\begin{align*}
 V(a,b,c)&=I(M(a,b),U(c,M(a,b))),\\
 P(a,b,c)&=I(c,M(c,b)),\\
 S(a,b,c)&=U\bigl(b,I(M(a,b),P(a,b,c))
                          \wedge I(M(c,b),c)\bigr),\\
 Q(a,b,c)&=U(c,I(S(a,b,c),N(P(a,b,c)))),\\
 G(a,b,c)&=N(I(V(a,b,c),Q(a,b,c)))\qquad(c\ne a,b).
\end{align*}
The remaining heads are the unrestricted diagonal
$\Gamma(B)=I(B,U(B,B))$ and the dominated head
\[
 \rho_8(D,A,H,u)=I(I(I(H,u),D),I(I(D,H),I(A,H))),\qquad u\le D,
\]
These are exactly the formulas of the inherited ordinary calculus.
Its only inference rules are modus ponens, from $A$ and $I(A,B)$ to $B$,
and theorem generalization, from $B$ to $U(i,B)$ for any label $i$.
In particular, $h_6$ has its original premise order. There are no
primitive equality, substitution, or semantic-truth rules.

Write $[a;w]=p(a,w)$, with total projections satisfying
\[
 \operatorname{hd}[a;w]=a,\quad\operatorname{tl}[a;w]=w,
 \qquad\operatorname{hd}(0)=\operatorname{tl}(0)=0.
\]
These are numerical stacks, so $[0;0]=0$ is allowed throughout.
For a store $(t,w,d,h)$ let $D=d+x$ and $H=h+y$. The command is
\begin{equation}
 \mathsf A_{c,x,y,z}(t,w,d,h)
       =(0,[q_0,q_1,\ldots,q_9,t;w],z,0),                 \tag{GD-A}
\end{equation}
where
\[
\begin{array}{c|l}
j&q_j\\\hline
0&G(H,D,c)\text{ if }c\ne H,D;\text{ otherwise }0\\
1&\Gamma(D)\\
2&h_6(D,c,H)\\
3&h_5(D,H)\\
4&h_4(D,c,H)\\
5&h_3(D,c,H)\\
6&h_2(D,c,H)\\
7&h_1(D,H)\\
8&\rho_8(D,c,H,d)\\
9&U(c,h).
\end{array}
\]
The source binding must establish this exact free-term equation, not
merely agreement on small numerical values. The fourth operand $z$ is
consumed by the actual fallthrough MP reader: after the push, its visible
inputs are zero and it compares $I(0,z)$ with zero. For $z=0$ the branch
succeeds and for $z>0$ it fails, but both branches leave top and cache
zero and offset $z$. Thus (GD-A) is asserted for every $z$.

The other commands are
\begin{equation}
 \mathsf S(t,w,d,h)=
 \begin{cases}
  \mathsf{halt},&h=1,\\
  (h,\operatorname{tl}w,d,\operatorname{hd}w),&h\ne1,
 \end{cases}                                               \tag{GD-S}
\end{equation}
and, writing $b=d+x$,
\begin{equation}
 \mathsf{MP}_x(t,w,d,h)=
 \begin{cases}
  (b,w,b,0),&I(t,b)=h,\\
  (0,w,b,0),&I(t,b)\ne h.
 \end{cases}                                               \tag{GD-M}
\end{equation}
Only Shift halts. Every $q_0,\ldots,q_8$ has residue zero modulo four;
$q_9$ has residue two and is at least two. No emitted $q_j$ is one.
These facts hold for all incoming stores, not only sound stores.

\subsection{Exact local words, including the early halt cases}

Words execute from left to right. Put $A_0=\mathsf A_{0,0,0,0}$ and
\[
 \mathsf{Push}=A_0;\mathsf S^{10};\mathsf{MP}_0,
 \qquad\mathsf{Pop}=A_0;\mathsf S^{12},
 \qquad\mathsf{Discard}=\mathsf{MP}_0.
\]
The first ten Shifts after an AX give
$(q_8,[t;w],0,q_9)$, and the first eleven give
$(q_9,w,0,t)$ when $z=0$. All those Shifts are nonhalting: their tests
see zero and then $q_0,\ldots,q_9$, never the saved top. Consequently
\begin{align*}
 \mathsf{Push}(t,w,d,h)&=(0,[t;w],0,0),\\
 \mathsf{Pop}(t,[a;w],d,h)&=
 \begin{cases}
   \mathsf{halt},&t=1,\\
   (t,w,0,a),&t\ne1,
 \end{cases}\\
 \mathsf{Discard}(t,w,0,h)&=(0,w,0,0).
\end{align*}
Push's MP compares an implication code of residue zero against $q_9$
of residue two, so it fails and returns zero. Pop's twelfth Shift tests
the saved top and hence halts exactly when $t=1$. Discard has the same
zero output in both branches. These derivations use total projections
and include the case where several displayed zero stack entries collapse
to numerical zero.

For any emitted position $0\le j\le8$, begin from $(t,w,0,0)$ and use
\begin{equation}
 \mathsf{Discard};\mathsf A_{c,D,H,0};
             \mathsf S^{j+2};\mathsf{Pop}^{9-j}.            \tag{GD-Sel}
\end{equation}
After the initial Shifts the top is $q_j$, the cache is $q_{j+1}$,
and the residual entries above $w$ are $q_{j+2},\ldots,q_9,0$.
There are exactly $9-j$ such entries. Each Pop preserves the non-one
top $q_j$, removes one entry, and places it in the cache. The last
removed entry is the saved zero. Thus the result is $(q_j,w,0,0)$.

Every ordinary axiom instance is selectable. For a gate $G(a,b,c)$,
choose AX operands $(c,b,a,0)$ and $j=0$. For an arbitrary diagonal
$\Gamma(B)$, use the specifically changed selector
\begin{equation}
 \mathsf{Discard};\mathsf A_{0,B,0,0};
                       \mathsf S^3;\mathsf{Pop}^8.         \tag{GD-Diag}
\end{equation}
For heads $h_1$ through $h_6$, use their positions $7,6,5,4,3,2$
in (GD-Sel) and set $(c,D,H)$ to the requested parameters. For the
restricted propositional leaf with $u\le D$, define
$\mathsf{Offset}_u=\mathsf{MP}_u$ and use
\begin{equation}
 \mathsf{Discard};\mathsf{Offset}_u;
 \mathsf A_{A,D\mathbin{\dot-}u,H,0};
                        \mathsf S^{10};\mathsf{Pop}.       \tag{GD-Rho}
\end{equation}
Offset sends $(0,w,0,0)$ to $(0,w,u,0)$; the selected accumulated
offset is $u+(D\mathbin{\dot-}u)=D$. The fourth AX operand clears
the outgoing offset. Its selected top is the requested $\rho_8$ and
is not one, so Pop is safe.

The new generalization word is shorter than the inherited AX11H word:
\begin{equation}
 \mathsf{Gen}_i=\mathsf{Push};\mathsf S;
                       \mathsf A_{i,0,0,0};\mathsf S^{11}. \tag{GD-Gen}
\end{equation}
Starting from $(a,w,0,0)$, Push gives $(0,[a;w],0,0)$ and the next
Shift gives $(0,w,0,a)$. Its halt test sees zero. AX then emits
$q_9=U(i,a)$ and saves top zero. The eleven-Shift identity returns
$(U(i,a),w,0,0)$. In particular $a=1$ is permitted: the AX clears
that cache before another Shift can test it. An Offset instruction is
neither needed nor silently assumed in (GD-Gen).

Finally, from $(0,[a,I(a,b),0;w],0,0)$ the suffix
\[
          \mathsf S;\mathsf S;\mathsf{MP}_b;\mathsf{Pop}
\]
halts if $a=1$ or $b=1$, and otherwise returns $(b,w,0,0)$.
The first two Shifts make $a$ the top and $I(a,b)$ the cache, unless
the second halts at $a=1$. MP then succeeds. The reserved zero under
the premises allows Pop to clear the offset and restore $w$ without
consuming any part of that base stack. This proof covers $a=b=w=0$.

\subsection{A primitive recursive compiler and its weak-base induction}

Encode a numerical derivation as a finite list with earlier-line pointers.
An axiom line is compiled by (GD-Sel), (GD-Diag), or (GD-Rho).
A generalization line appends (GD-Gen) to its premise word.
If $e_1$ proves $I(a,b)$ and $e_0$ proves $a$, its MP word is
\[
 \mathsf{Discard};\mathsf{Push};W_{e_1};\mathsf{Push};
 W_{e_0};\mathsf{Push};\mathsf S;\mathsf S;
                         \mathsf{MP}_b;\mathsf{Pop}.
\]
The first Discard and Push reserve the zero needed by the final suffix.
For every derivation $e$ the exact contract is
\begin{equation}
 W_e(t,w,0,0)=\mathsf{halt}\quad\text{or}\quad
 W_e(t,w,0,0)=(\operatorname{concl}(e),w,0,0).              \tag{GD-Cmp}
\end{equation}
Earlier halts remain absorbing and require no special inference about
the conclusion of the subproof.

Here is an explicit quantifier-free induction supporting this usual
proof-by-compilation argument. Let the list have $N>0$ lines; compile
all their words by a bottom-up primitive recursion. Let $L$ be their
maximum length and $K$ their maximum command code. For a primitive
recursive encoded command step define the monotone majorant
\[
 \Phi_K(s)=1+\max\left(s,
     \max_{c\le K,\,v\le s}\operatorname{next}(c,v)\right).
\]
It is primitive recursive, and a prefix of at most $L$ commands started
at a store code at most $s$ stays below $\Phi_K^L(s)$.
Fix an arbitrary normalized input code $s_0$, and put
\[
 B(j)=\Phi_K^{L(N-1-j)}(s_0)\quad(j<N).
\]
For $i<j$, $\Phi_K^L(B(j))\le B(i)$. Induct on the bounded conjunction
of (GD-Cmp) for every earlier line $i<j$ and every normalized input
code at most $B(i)$. Word evaluation and this bounded conjunction have
primitive recursive characteristic functions. The induction is therefore
quantifier free. Every premise call in a word for line $j$ is preceded
by at most $L$ commands, so its actual input lies within the corresponding
earlier bound. The local contracts prove the induction step, including
all early halts. At the final line $B(N-1)=s_0$, giving (GD-Cmp) for
the arbitrary input. This does not assume induction on an unbounded
assertion about every base stack.

If $e$ concludes one, the word
\begin{equation}
                       V_e=A_0;W_e;\mathsf{Pop}             \tag{GD-Ref}
\end{equation}
halts from every incoming store. AX first normalizes offset and cache,
leaving a numerical base stack to which (GD-Cmp) applies. If no earlier
halt occurred, the result has top one and cache zero; Pop detects that
one. Merely appending a Shift would not suffice with a cached-only test.

\subsection{Fair scheduling and syntactic soundness}

The same four-field decoder is used. Define
$e(\epsilon)=0$ and $e(c::v)=1+b_c(e(v))$, where
\[
 b_{\mathsf S}(z)=p(0,z),\qquad
 b_{\mathsf{MP}_x}(z)=p(1,p(x,z)),
\]
\[
 b_{\mathsf A_{c,x,y,u}}(z)=p(c+2,p(y,p(x,p(u,z)))).
\]
With fresh counter $n$ and continuation $\ell$, the scheduler parses
raw seed $n$ and advances the counter when $\ell=0$; otherwise it
parses $\ell-1$ and retains the counter. Each remaining continuation
is a nested Cantor component of its raw seed, so a positive continuation
strictly decreases. A fresh candidate $n$ therefore drains in at most
$n+1$ commands unless a halt occurs earlier.

Starting from fresh counter two, prefix another AX before (GD-Ref).
Its raw candidate
\[
 N_e=p(2,p(0,p(0,p(0,e(V_e)))))\ge3
\]
executes that AX and then the all-store halting word, regardless of the
store carried by earlier candidates. A loose explicit command-clock
bound is $(N_e+2)^2$, covering all earlier drains and the successful
candidate. The encoding, decrease argument and finite quadratic bound
are proved by quantifier-free induction. No reset between trials or
unbounded search for a first successful candidate is assumed.

For soundness annotate the top, cache and canonical Cantor-stack entries
by finite $\mathcal C_4$ derivations. The initial zeros have the one-row
proof $\rho_8(0,0,0,0)=0$. AX preserves the annotation: $q_1=\Gamma(D)$
is an unrestricted diagonal axiom, $q_9=U(c,h)$ is generalization of
the cache proof at the explicit label $c$, and the actual restriction
on $q_8$ is $d\le d+x$. Other entries are the selected ordinary heads,
the guarded gate, or zero. Arbitrary outgoing $z$ imposes no proof
obligation because the offset register is not an asserted formula.
Successful MP combines the annotated premises; failure returns zero.
Nonhalting Shift transports annotations, using the zero proof for an
empty numerical stack. A halting Shift exposes an annotated cache equal
to one, hence a numerical refutation.

For $w>0$ its Cantor tail is strictly smaller than $w$, so the entire
canonical stack annotation can be checked by a bounded scan of at most
$w$ iterations. Annotation updates and their validity tests are primitive
recursive. Quantifier-free induction on the number of commands proves
soundness uniformly. Together with completeness this establishes
\[
 \mathsf{iPRA}\vdash
   \operatorname{Halt}(R_{\rm AX11GDH})
        \leftrightarrow\exists e\,\operatorname{Prf}_{\mathcal C_4}(e,1).
\]

\subsection{The literal-machine and PA endpoints}

The independent source binding expands all 3,671 primitive instructions
from the handwritten clear, transfer, duplicate, pairing, unpairing,
comparison and multiplication templates of the inherited
\texttt{source\_command\_completion.tex}. It compares every register
argument and outgoing edge, re-expands the named syntax, and independently
resolves that syntax's control tree. The resulting ports are fetch 25,
decode 34, AX 50, MP 3572, MP comparison 3642, MP success 3652,
return 3659 and Shift 3660. The primitive template equations and exact
clock functions are unchanged. They are proved by quantifier-free counter
induction in the inherited source-command proof; the finite duration
examples in the new checker supplement those uniform arguments.

The new free-term AX check substitutes those template equations in the
exact 190-token order. It retains both zero-on-equality guard operations,
checks the scratch precondition at every macro boundary, and compares
all nine endpoint registers. Starting from
\[
 (N,T,W,A,D,H,X,Y,L)
   =(n,t,w,c,d,h,0,0,p(y,p(x,p(z,\ell)))),
\]
the endpoint at port 3572 is
\[
 (n,0,[q_0,\ldots,q_9,t;w],0,0,0,0,0,p(z,\ell)).
\]
This is a formal arithmetic term identity, with $q_j$ as defined above,
and includes the saved raw cache in $q_9=U(c,h)$. The consumed offset
and cache in $q_1,\ldots,q_8$ are $D=d+x$ and $H=h+y$ as required.

For completeness the remaining matched source fragments have the
following universal consequences. At an ordinary MP entry, ReadDCons
replaces $(T,D,L)=(t,d,p(x,\ell))$ by
$(I(t,d+x),d+x,\ell)$ and preserves the stack and cache. Transfer places
that implication in $X$ and clears $T$. The destructive comparison clears
$H,X$; on equality the copy-and-restore template puts $d+x$ in $T$,
and on inequality it leaves $T=0$. This is precisely (GD-M), with
$A=X=Y=0$ at return. Applying the same argument to the displayed AX
endpoint gives $I(0,z)$ as the compared value. Pairing injectivity proves
$I(0,z)=0$ exactly when $z=0$; in that branch the copied value is zero,
and in the other branch the already-zero top is preserved. Both branches
therefore give (GD-A) for every fourth operand, with residual $\ell$.

For Shift, clearing $T$ and copying-and-restoring $H$ makes $T=h$;
incrementing the zero scratch register makes $X=1$. The root comparison
halts precisely when $h=1$. Otherwise it clears $H,X$, unpairs $W$, and
transfers its first component to $H$, giving (GD-S). Thus the test is
cached-only, including when the saved top is one. Finally, the fetch
prefix maps a positive residual $\ell$ to $\ell-1$ in two primitive
steps, or maps $(N,L)=(n,0)$ to $(n+1,n)$ in $7n+4$ steps. Unpairing
and the two dispatch decrements select exactly the scheduler's three
tags and leave $A=0$ on both non-AX branches. The AX guard clears $A$
on the third branch. All ordinary returns have the required zero scratch
registers and fetch port 25.

These statements supply a positive primitive recursive exact duration
for one complete scheduler command: add the fetch and dispatch costs to
the exact template costs evaluated at their successive arithmetic stores.
The fixed AX word has only 190 tokens, so this composition is primitive
recursive even when its values are extremely large. The corresponding
iterated primitive-step equality follows by the template theorems and
finite composition, not by the 882 local and 324 composed duration
examples. The resulting scheduler-to-primitive-source clock can be
composed with the separate physical source-instruction clocks.

The symmetric-Shift sibling AX11GDHS has exactly the same abstract
commands and AX body. Its independently checked source has 3,668
instructions; every named port through the return at 3659 is unchanged,
and Shift still starts at 3660. The eight-instruction Shift suffix is:
clear $T$; duplicate $H$ into $T,X$ while clearing $H$; increment $H$;
jump to the root comparison. Thus its comparison input is
$(T,H,X)=(h,1,h)$, instead of $(h,h,1)$. Equality is still exactly
$h=1$, and comparison clears $H,X$ in both cases. The common nonhalting
stack decoder therefore gives precisely (GD-S), including $h=0$ and
an arbitrary saved top. The clear and duplicate counter lemmas give an
exact positive primitive recursive clock for this suffix as well.
Its complete source binding independently matched all 3,668 instructions
and edges and the same AX free-term endpoint. Consequently every
compiler, scheduler, soundness and weak-base argument above applies
unchanged to AX11GDHS; its separate literal physical certificates supply
the corresponding source-instruction interface.

The literal machine is identified with this scheduler by the separate
exact free-term source binding, arbitrary-register physical contracts,
complete blank startup and primitive recursive positive-clock composition.
Those are essential hypotheses of this final composition, and bounded
numerical command tests alone do not replace them. With those bindings,
the new table and the ordinary AX11H parent halt under exactly the same
circumstances, by the two displayed scheduler-to-$\mathcal C_4$ theorems.
The same negative equivalence follows intuitionistically.

Composing with the explicitly cited updated AX11H written calibration
then gives
\[
 \mathsf{iPRA}\vdash
  \operatorname{Con}(\mathrm{PA})
        \leftrightarrow\neg\operatorname{Halt}_{\rm blank}(M_{\rm AX11GDH}).
\]
The new proof preserves the numerical calculus and PA predicate rather
than assuming that an unrelated short verifier has the same consistency
strength. Its absolute PA calibration is an inherited written theorem;
no new complete Lean or Rocq kernel derivation is asserted.

